% Mid-year project progress report for CS 510 Intro to Privacy
% Steve Willoughby & Jacob Olinger
% 
% VLDB template version of 2020-08-03 enhances the ACM template, version 1.7.0:
% https://www.acm.org/publications/proceedings-template
% The ACM Latex guide provides further information about the ACM template
\documentclass[sigconf, nonacm]{acmart}
\bibliographystyle{ACM-Reference-Format}
\usepackage{xcolor}
\usepackage[normalem]{ulem}
\definecolor{ontrack}{HTML}{4A6741}
\definecolor{behind}{HTML}{8F1402}
\definecolor{atrisk}{HTML}{888800}
\newcommand\tuple[1]{\langle#1\rangle}
\newcommand\var[1]{$\langle$\textit{#1}$\rangle$}

%% The following content must be adapted for the final version
% paper-specific
\newcommand\vldbdoi{XX.XX/XXX.XX}
\newcommand\vldbpages{XXX-XXX}
% issue-specific
\newcommand\vldbvolume{14}
\newcommand\vldbissue{1}
\newcommand\vldbyear{2020}
% should be fine as it is
\newcommand\vldbauthors{\authors}
\newcommand\vldbtitle{\shorttitle} 
% leave empty if no availability url should be set
\newcommand\vldbavailabilityurl{URL_TO_YOUR_ARTIFACTS}
% whether page numbers should be shown or not, use 'plain' for review versions, 'empty' for camera ready
\newcommand\vldbpagestyle{plain} 

\begin{document}
\title{IoT Access Control}
\subtitle{Final Project Report \\ June 2025}

%%
%% The "author" command and its associated commands are used to define the authors and their affiliations.
\author{Steve Willoughby}
\affiliation{%
  \institution{Portland State University}
  \streetaddress{P.O. Box 751}
  \city{Portland}
  \state{Oregon}
  \country{USA}
  \postcode{97207-0751}
}
\email{slw3@pdx.edu}
\author{Jacob Olinger}
\affiliation{%
  \institution{Portland State University}
  \streetaddress{P.O. Box 751}
  \city{Portland}
  \state{Oregon}
  \country{USA}
  \postcode{97207-0751}
}
\email{jacoboli@pdx.edu}
\maketitle
\section{Abstract}
Internet of Things (IoT) systems include a multitude of sensors which ``publish''
data streams to another multitude of ``subscribers''. Between them sits a ``broker''
which routes these published messages to all the subscribers interested in them.
However, there is a need for more robust access control to govern who can
subscribe to potentially sensitive data, and under what conditions they may do so.
Complicating things further, even an authorized subscriber may only have a valid
need to access sensor data at certain times, in certain locations or under certain
circumstances.

\section{Introduction}
Much of our increasingly-automated infrastructure relies on small, internet-connected devices, or the so-called ``Industrial Internet of Things'' (IIoT).\footnote{As this class project is, with instructor approval, a portion
of the research work toward Steve's Ph.D. dissertation, the text in the introduction section and other areas in this report is quoted from his research work in progress on this topic without further explicit reference, as this is considered all part of the same research work being done, and the results of this project will eventually be incorporated into the published papers and dissertation work as well.}
Among the numerous protocols used to enable these technologies to function is the MQ Telemetry Transport protocol (MQTT). This is a simple,
lightweight plaintext message passing protocol which is designed to put a low overhead on small, low-power devices. As a result, it does not by default
support encryption, authentication, message integrity validation, or access controls. It is an example of a ``publisher/subscriber'' system 
(commonly ``pub/sub'' for short) where sensor devices collect readings and other telemetry data as ``topic publishers'', sending them to a central broker service. This, in turn, retransmits the telemetry data to client systems, other servers, or even other IoT devices such as device controllers which
``subscribe'' to that data. Further, other independent agents which are in control of, e.g., automated building systems, publish messages to devices to adjust
building airflow or activate door locks or whatever other device tasks they are in charge of, all through MQTT messages moving through the system.

\section{Related Work}
The quest to effectively and efficiently manage subscriber access to published data (as well as the authority for publishers to provide that
data) is not new, and a number of approaches have been proposed to date. In their 2023 paper \cite{ma2022tbac}, Ma and Yang describe a topic
oriented model which grants access to certain subscribers according to which topics they are allowed to see.  Additionally, Xie, Shi, et al. pointed
out \cite{xie2020topic} another important consideration for managing such access controls---dynamic rulesets. Their controls were based on the idea
of mapping a set of users, devices, and topics via rules as to which actors can publish or subscribe to which topics.

In their previous work \cite{colombo2018access}, Colombo \& Ferrari proposed a design for an access control manager proxy which can be deployed into an MQTT message broker to enforce
a set of policies, possibly even causing messages to be rerouted to specific subscribers, held, or dropped entirely per those policy decisions.


\section{Project Description}
Our project seeks to add fine-grained, high-performance access control to MQTT pub/sub message broker systems for IoT networks.

It differs from the current state of the art in that most current offerings only go as far as basic ACL rulesets.
We seek to build in particular on the aforementioned prior work such as
that of Colombo \& Ferrari, although we will more tightly integrate our
policy enforcement engine with the broker to optimize performance as well as try to achieve a more fine-grained ABAC policy model to
control subscriber access to published message topics.

\subsection{Access Control Policy Model}
\sout{Access to the published messages is, due to the topic-oriented nature of MQTT IoT traffic, is essentially a fine-grained ABAC policy model where the MQTT topic hierarchy is the primary subject identifier to select the resource being requested by the subscribers.}

\sout{They are allowed to subscribe to a specific message topic or to a ``wildcard pattern'' which may match a number of topics as they are published and arrive at the broker for distribution to the subscribing clients, which presents some challenges for our policy enforcement engine to take into consideration when resolving our access control rules.}

\sout{We represent access to a given message as a tuple }
$$\tuple{t,V,P,m,po}$$
\sout{where $t$ is the MQTT topic pattern under which the message is published via the MQTT broker. Theoretically,
this could map to an arbitrary kind of resource in the real world but as far as we are concerned and the
MQTT protocol will see it, these are all ``messages'' being published, routed, and sent to subscribers, so we will
refer to them as such henceforth. MQTT topics are hierarchical in structure, with slashes separating
each level in the hierarchy, and special wildcards ``{\tt +}'' and ``{\tt \#}'' may be used to form patterns to match
many topics in a single rule. For example, ``{\tt bldg1/+/temp/\#}'' might be part of a tuple to specify a rule that controls
access to messages about temperature readings from all thermostats on all floors of building 1.}

\sout{The next value, $V$, specifies a set of times during which
the tuple is considered to be valid, i.e., in effect. $P$ lists a set of access control policies to be enforced predicated on the conditions of the other elements
of the tuple.}\footnote{\sout{Having $V$ separate from the policy set $P$ is an optimization step, since a range of times throughout the day is a common use case for IoT systems, and this allows us to set up our approved subscription topic tree to include approved access times which we can continue to use when routing messages rather than have those times cause policy rechecks throughout the day just for their sake.}}

\sout{The value $m$ specifies contextual information for the rule that is dependent on content of the message payload data itself.}

\sout{Finally, the $po$ value names the policy owner in order to document ownership of access management data and for maintenance and audit purposes.}

\sout{Thus, anyone who wishes to subscribe to receive all messages which are published to topics matching the pattern $t$, for all valid time values $v\in V$, must satisfy all the conditions of all access policies $p\in P$ including any per-message policies indicated by $m$ (if any).}

\subsection{Policy Model}
\sout{In our concept of published topics \emph{vis-\`a-vis} access control, we 
consider the role of the MQTT broker to be to route messages from publishers to subscribers where those messages are in the intersection of the set of topics that match those that each subscriber has requested to receive, and the set of valid times and the set of policies that grant access to the subscriber for those topics.}

\sout{Expanding further on the definition of the set $P$ of policies in the tuple that defines access to a topic, each element of $P$ is itself a tuple of values} $$\tuple{e,c,g,o}$$ \sout{where $e$ is an expression including the subject's ABAC attributes relevant to determining the access decision, $c$ is the context in which the subject is operating (i.e.
the purpose for the access), $g$ is the resulting grant in the event
the expression is satisfied, and $o$ is the operation for which this
applies.}

\sout{The expression $e$ may include the ABAC attribute names. By default, the name alone has a {\tt true} or {\tt false} value if it has boolean type and {\tt true} or {\tt false} value, respectively; the prefix operator {\tt !} negates its boolean value. Thus simple attribute expressions such as ``{\tt MANAGER}'' or ``{\tt !GUEST}'' are easily specified. Logical operators {\tt\&} and {\tt|} may be employed for multiple attributes allowing expressions such as ``{\tt TEAM1|TEAM2}''.}

\sout{For attributes with non-boolean values, standard comparison operators 
{\tt<},
{\tt<=},
{\tt>},
{\tt>=},
{\tt=},
{\tt!=},
and numeric and string constant values may be used along with parentheses for grouping (see Fig.~\ref{abac-ops}).}
\begin{figure}
    \begin{center}
    \begin{tabular}{llll}
    $b$&bool value&
    $x$\tt=$y$&$x=y$\\
    \tt!$b$&not bool value&
    $x$\tt!=$y$&$x\ne y$\\
    $x$\tt<$y$&$x<y$&
    $x$\tt<=$y$&$x\le y$\\
    $x$\tt>$y$&$x>y$&
    $x$\tt>=$y$&$x\ge y$\\
    $x$\tt\&$y$&$x\wedge y$&
    {\tt(}\dots{\tt)}&grouping\\
    $x$\tt|$y$&$x\vee y$\\
    \end{tabular}
    \caption{ABAC Attribute Expression Operators\label{abac-ops}}
    \end{center}
\end{figure}

\sout{An empty value for $e$ is considered {\tt true} as a special case. If $e$ is {\tt true}, the rule is processed. Otherwise, it is skipped and the next rule in the set is attempted, in order.}

\sout{The next value in the tuple, $g$, specifies the grant and may be either {\tt ALLOW} or {\tt DENY}. This terminates the
evaluation of the rule set.}

\sout{Finally, the $o$ value specifies the operation which is to be granted; this is always {\tt SUBSCRIBE} for now, but may include more operations in our future work.}

\sout{If the end of the policy list is reached without any of the attribute expressions $e$ evaluating to {\tt true}, then
{\tt DENY} is assumed by default. This can be amended by adding a catch-all rule that matches all topics for all time ranges for all subjects unconditionally and grants approval.}

\subsection{Scope of Investigation}
\sout{For the purposes of our project, we will consider the following set of access rules to be imposed on MQTT traffic and measure the overhead of applying the enforcement of those rules in our broker.}
{\color{red}\dots}

\section{Timeline}
Our original estimated timeline was:
\begin{description}
	\item[\color{ontrack}15 Apr] Proposal Presentation.
	\item[\color{ontrack}19 Apr] Project scope and plan finalized following feedback.
	\item[\color{ontrack}26 Apr] Design for experimental simulated IoT network.
	\item[\color{ontrack}03 May] Implementation of IoT clients and brokers.
	\item[\color{behind}10 May] Control runs completed, Mid-Quarter Report.
	\item[\color{atrisk}17 May] Add access controls to broker, 2nd experiments.
	\item[\color{atrisk}24 May] Additional experiments and scale tests.
	\item[31 May] Complete data analysis.
	\item[05 Jun] Final Report ready.
\end{description}
\subsection{Progress to Date}
\sout{As suggested by the color-coded dates above, we are about a week behind schedule due to the time
required to search for suitable IoT client programs to use for publishers and subscribers and for
the time to fully settle on the access control policy model to be used in the broker modifications
we will be employing.}

\sout{We have obtained the clients and have an installed broker instance, and are just about ready to begin the
control runs now.}

\section{Results and Analysis}
\sout{We will set up simulated traffic and measure the performance impact of our new
controls on the base system}

\sout{We will run a number of test cases to demonstrate that the access controls are effective
at limiting access to only authorized subscribers who meet the policy restrictions within
the given domains.}

\subsection{Experimental Setup}
We are working in a compartmentalized environment, which allows us to develop each of our areas of concern---the server (MQTT broker) modifications and the clients (simulated publishing and subscribing devices)---in relative isolation. The clients
are in Docker containers, allowing them to be instantiated identically as many times as needed in the experimental environment, each running inside its isolated container but able to send and receive network traffic with the server.

This setup allows each team member to separately develop the server and client components away from the actual test environment as well, allowing development and debugging in a way that does not interfere with running experiments.

\subsection{Orchestration}
\sout{The operation of the experimental lab is governed by two main scripts.}

\sout{The launch script will start up the server (MQTT broker) and then the publishers and subscribers according to its configuration.}

\sout{We then collect the data after the run is complete and analyze the data.}

\subsubsection{Configuration}
\sout{The experimental runs are controlled by a JSON file which governs how the clients and servers are configured as well as what the experimental parameters will be. The format of this file and the meaning of these parameters are explained by example below:}

\sout{The clients write stats about the data they receive to data collection files in the \var{collectiondir}{\tt/}\var{jobid}{\tt/}\var{rawdata} directory under a unique file for each running client instance. The launch script tracks the running clients and manages
which ones are still running until they're all done and we have all their data.}

\subsubsection{Data Collection}
\sout{The containerized clients receive their subscribed data and record their results into a data file with one line per received message in the form:}

\begin{center}
\var{datestamp} \var{topic} \var{bytes\_received}
\end{center}

\sout{The data collection script comes along after the experimental runs are completed and scans all these data files, tabulating the data and
analyzing the results to show the traffic speed through the broker,
which we can compare to the speed of the same traffic with and without
access controls applied. We will also use this to verify that messages
were or were not filtered as appropriate per the access control policies.}

\subsubsection{Extensibility and Updates}
\sout{This is a baseline to get started. As we build out the capabilities of the experimental system we'll add more parameters related to access control to how we launch the clients and server for the experiments we'll run.}

\subsection{Hardware Profile}
The server used to perform all the experimental runs, and upon which our performance benchmarks is based, was an Intel\textregistered\ Core\texttrademark\ i9-10980XE (36 cores) running at 4.80\,GHz configured with 128\,Gb of memory.
It was running the Arch-based Garuda Linux distribution 6.13.7 Zen Linux kernel at the time of the experimental results reported here.
\section{Conclusions}
\sout{Pending results from experimental work.}
\section{Individual Contributions}
The efforts were divided between Steve and Jacob along the boundaries of broker and client systems,
with Steve handling the broker, including the setup of the Mochi MQTT broker and making the access
control enforcement modifications to it, and Jacob setting up the containerized instances of
the MQTT client programs acting as publishers and subscribers as well as the orchestration of
those clients' actions in sending their traffic as part of our experimental runs.

%
\bibliography{final-willoughby-olinger}
\end{document}
% 2-4 pp
%Introduction
%Describe the problem that you are addressing in this project and why it is interesting in the context of Privacy-aware Computing. Include examples or case studies to motivate the problem. Briefly describe what you hope to achieve at the completion of the project.
%	
%This criterion is linked to a Learning Outcome Related Work
%Overview of previous work related to project including research papers, projects, applications, or regulations. Importantly, how does your project relate to these existing approaches?
%
%Cite at least 3 other works.
%	
%This criterion is linked to a Learning Outcome Project Description
%A detailed description of the project including details of the specific privacy problem that you are addressing, definitions of terminology, details of the solution/approach, how the solution will improve privacy or privacy-awareness, how you are implementing the solution so far - details of tools/algorithms/frameworks that you are using, how you plan to evaluate your approach and what will be the metrics used. Include any architectural diagram or system diagrams of your approach. This will most likely be largest section of the report.
%	
%This criterion is linked to a Learning Outcome Timeline
%Include details of the tasks completed in the project and milestones achieved as well as the remaining tasks and how you plan to complete before the project presentation on March 10th. Optionally, include mention any technical or logistical challenges your team is facing and how you plan to address them.
%	
%This criterion is linked to a Learning Outcome Others
%Project title + Abstract
%Individual Contributions
%Results & Analysis
%This criterion is linked to a Learning Outcome References and Writing
%Correctly cited references
%Correctness of writing (Grammar, readability, spelling check)
